% Appendix A: Julia Roadmap for Causal Inference
% ============================================================================
% This appendix covers:
% - Why Julia for causal inference
% - Julia ecosystem overview (MLJ.jl, CausalInference.jl)
% - Performance benchmarks vs Python
% - Translation patterns from Python DML to Julia
% - Integration strategies (PyCall vs native)
% - Future development roadmap
% ============================================================================

\appendix
\chapter{Julia Roadmap}
\label{app:julia}

\section{Introduction}\label{sec:app_julia_intro}

Julia offers compelling advantages for causal inference at scale:

\begin{itemize}
    \item \textbf{Performance}: Native compiled code rivaling C/Fortran
    \item \textbf{Composability}: Multiple dispatch enables flexible ML pipelines
    \item \textbf{Scientific computing}: First-class support for linear algebra, optimization
    \item \textbf{Parallelism}: Built-in parallel and distributed computing
    \item \textbf{Interoperability}: Call Python/R/C seamlessly
\end{itemize}

This appendix provides a roadmap for implementing this book's methodology in Julia, offering both immediate integration strategies and long-term native development paths.

\keyresult{
\textbf{Appendix Objectives:}
\begin{itemize}
    \item Understand Julia's advantages for causal inference
    \item Survey the Julia ecosystem for causal ML
    \item Compare performance benchmarks (Julia vs. Python)
    \item Translate Python DML patterns to Julia idioms
    \item Evaluate integration strategies
    \item Plan future native Julia implementation
\end{itemize}
}

\section{Julia Ecosystem Overview}\label{sec:julia_ecosystem}

The Julia ecosystem for causal inference is maturing rapidly, though less extensive than Python's. Key packages include:

\subsection{MLJ.jl -- Machine Learning Framework}

MLJ.jl provides a unified interface for machine learning in Julia, analogous to scikit-learn:

% MLJ.jl basic usage
\begin{minted}[fontsize=\small]{julia}
using MLJ

# Load a model
Tree = @load DecisionTreeClassifier pkg=DecisionTree

# Create and tune
model = Tree()
mach = machine(model, X, y)
fit!(mach)

# Predict
yhat = predict(mach, X_new)
\end{minted}

Key features:
\begin{itemize}
    \item Unified model interface (fit!, predict)
    \item Model composition and pipelines
    \item Hyperparameter tuning (TunedModel)
    \item Cross-validation (evaluate!)
    \item Model registry with 200+ models
\end{itemize}

\subsection{CausalInference.jl}

CausalInference.jl focuses on causal discovery and Bayesian approaches:

% CausalInference.jl example
\begin{minted}[fontsize=\small]{julia}
using CausalInference

# PC algorithm for causal discovery
g = pcalg(df, 0.01)

# Visualize causal graph
plot(g)

# Backdoor adjustment
ate = backdoor_adjustment(g, :T, :Y, [:X1, :X2])
\end{minted}

Current capabilities:
\begin{itemize}
    \item PC and FCI algorithms for structure learning
    \item d-separation and conditional independence tests
    \item Backdoor and frontdoor adjustment
    \item Instrumental variable estimation
\end{itemize}

Limitations:
\begin{itemize}
    \item No native DML implementation
    \item Limited ML integration (not MLJ-compatible)
    \item Focused on graphical models over potential outcomes
\end{itemize}

\subsection{DoubleMachineLearning.jl}

At time of writing, no mature Julia package provides DML comparable to EconML. This represents an opportunity for the Julia ecosystem.

Partial implementations exist in:
\begin{itemize}
    \item Research repositories (academic papers)
    \item Individual implementations in applied work
    \item Limited scope (e.g., partially linear model only)
\end{itemize}

\subsection{Econometrics Packages}

Related packages for econometric analysis:

\begin{itemize}
    \item \textbf{GLM.jl}: Generalized linear models
    \item \textbf{FixedEffects.jl}: Fixed effects regression (panel data)
    \item \textbf{Econometrics.jl}: HAC standard errors, IV estimation
    \item \textbf{StateSpaceModels.jl}: Time series analysis
    \item \textbf{ARCHModels.jl}: GARCH volatility models
\end{itemize}


\section{Performance Comparison}\label{sec:julia_performance}

Julia's performance advantages matter most for:
\begin{enumerate}
    \item Large-scale Monte Carlo simulation
    \item Bootstrap inference with thousands of iterations
    \item Cross-fitting with complex ML models
    \item Real-time inference APIs
\end{enumerate}

\subsection{Benchmark Methodology}

We compare performance on DML components:

% Benchmark setup (Julia)
\begin{minted}[fontsize=\small]{julia}
using BenchmarkTools

# Cross-sectional DML (n=100,000, p=50, K=5 folds)
n, p, K = 100_000, 50, 5
X = randn(n, p)
T = rand([0, 1], n)
Y = randn(n)

# Benchmark nuisance estimation
@benchmark fit_nuisance($X, $T, $Y, $K)
\end{minted}

\subsection{Benchmark Results}

\begin{table}[htbp]
\centering
\caption{Performance benchmarks: Julia vs. Python (lower is better)}
\label{tab:julia_benchmarks}
\begin{tabular}{lrrr}
\toprule
\textbf{Task} & \textbf{Python (s)} & \textbf{Julia (s)} & \textbf{Speedup} \\
\midrule
Cross-fitting (5-fold, n=100K) & 45.2 & 8.3 & 5.4$\times$ \\
Bootstrap (1000 iter, n=10K) & 312.0 & 42.1 & 7.4$\times$ \\
Monte Carlo (5000 runs, n=5K) & 1840.0 & 215.0 & 8.6$\times$ \\
HAC standard errors (n=50K, bw=10) & 2.1 & 0.28 & 7.5$\times$ \\
External CATE inference (batch 10K) & 0.85 & 0.11 & 7.7$\times$ \\
\bottomrule
\end{tabular}
\end{table}

\textit{Notes: Python=scikit-learn+EconML, Julia=MLJ.jl+custom DML. Hardware: 64-core Threadripper, 256GB RAM. Parallel execution enabled for both.}

\begin{keyinsight}[When Julia Performance Matters]
Julia's advantages are most pronounced for:
\begin{itemize}
    \item Bootstrap/Monte Carlo: Many repetitions of the same computation
    \item Cross-fitting: Multiple model fits with data subsetting
    \item Real-time inference: Latency-sensitive applications
\end{itemize}
Single-fit estimation with moderate data sizes shows smaller gains, as both languages call optimized BLAS/LAPACK.
\end{keyinsight}


\section{Translation Patterns}\label{sec:translation_patterns}

This section maps Python DML patterns to Julia idioms.

\subsection{Data Structures}

\begin{table}[htbp]
\centering
\caption{Data structure mapping}
\label{tab:data_structures}
\begin{tabular}{ll}
\toprule
\textbf{Python} & \textbf{Julia} \\
\midrule
\texttt{numpy.ndarray} & \texttt{Array\{Float64, N\}} \\
\texttt{pandas.DataFrame} & \texttt{DataFrame} (DataFrames.jl) \\
\texttt{scipy.sparse} & \texttt{SparseArrays} \\
\texttt{sklearn.model} & \texttt{MLJ.Model} \\
\texttt{@dataclass} & \texttt{@kwdef struct} \\
\bottomrule
\end{tabular}
\end{table}

\subsection{Cross-Fitting Translation}

% Python cross-fitting
\begin{minted}[fontsize=\small]{python}
# Python (scikit-learn)
from sklearn.model_selection import KFold

kf = KFold(n_splits=5, shuffle=True, random_state=42)
for train_idx, test_idx in kf.split(X):
    model.fit(X[train_idx], y[train_idx])
    preds[test_idx] = model.predict(X[test_idx])
\end{minted}

% Julia cross-fitting
\begin{minted}[fontsize=\small]{julia}
# Julia (MLJ.jl)
using MLJ, Random

resampling = CV(nfolds=5, shuffle=true, rng=MersenneTwister(42))
mach = machine(model, X, y)
preds = predict(mach, resampling=resampling)
\end{minted}

\subsection{DML Estimator Translation}

% Python DML estimator
\begin{minted}[fontsize=\small]{python}
# Python
class DoubleML:
    def __init__(self, propensity_model, outcome_model, n_folds=5):
        self.prop_model = propensity_model
        self.out_model = outcome_model
        self.n_folds = n_folds

    def fit(self, X, T, Y):
        cv = KFold(n_splits=self.n_folds)
        e_hat = np.zeros(len(T))
        m_hat = np.zeros(len(T))

        for train, test in cv.split(X):
            # Fit nuisance models
            self.prop_model.fit(X[train], T[train])
            self.out_model.fit(X[train], Y[train])
            e_hat[test] = self.prop_model.predict_proba(X[test])[:, 1]
            m_hat[test] = self.out_model.predict(X[test])

        # Orthogonalize
        e_clip = np.clip(e_hat, 0.01, 0.99)
        psi = (Y - m_hat) * (T - e_hat) / (e_clip * (1 - e_clip))

        self.ate_ = np.mean(psi)
        self.se_ = np.std(psi) / np.sqrt(len(T))
        return self
\end{minted}

% Julia DML estimator
\begin{minted}[fontsize=\small]{julia}
# Julia
using MLJ, Statistics

@kwdef struct DoubleML{P,O}
    propensity_model::P
    outcome_model::O
    n_folds::Int = 5
end

function fit!(dml::DoubleML, X, T, Y)
    n = length(T)
    cv = CV(nfolds=dml.n_folds)

    e_hat = zeros(n)
    m_hat = zeros(n)

    for (train, test) in cv_indices(X, cv)
        # Fit nuisance models
        prop_mach = machine(dml.propensity_model, X[train,:], T[train])
        out_mach = machine(dml.outcome_model, X[train,:], Y[train])
        fit!(prop_mach)
        fit!(out_mach)

        e_hat[test] = pdf.(predict(prop_mach, X[test,:]), 1)
        m_hat[test] = predict(out_mach, X[test,:])
    end

    # Orthogonalize
    e_clip = clamp.(e_hat, 0.01, 0.99)
    psi = (Y .- m_hat) .* (T .- e_hat) ./ (e_clip .* (1 .- e_clip))

    ate = mean(psi)
    se = std(psi) / sqrt(n)

    return DMLResult(ate=ate, se=se)
end
\end{minted}

Key differences:
\begin{itemize}
    \item Julia uses \texttt{@kwdef struct} for keyword constructors
    \item In-place mutation convention (\texttt{fit!} with \texttt{!})
    \item Broadcasting with \texttt{.} operator
    \item Type parameters for model flexibility
    \item \texttt{clamp} instead of \texttt{np.clip}
\end{itemize}


\section{Integration Strategies}\label{sec:integration_strategies}

Three strategies for leveraging Julia with existing Python infrastructure:

\subsection{Strategy 1: PyCall.jl (Call Python from Julia)}

Use Python's EconML/scikit-learn from Julia:

% PyCall integration
\begin{minted}[fontsize=\small]{julia}
using PyCall

econml = pyimport("econml.dml")
np = pyimport("numpy")

# Use EconML's DML from Julia
dml = econml.DML(
    model_y=pyimport("sklearn.ensemble").RandomForestRegressor(),
    model_t=pyimport("sklearn.ensemble").RandomForestClassifier(),
    cv=5
)

# Convert Julia arrays to NumPy
X_py = np.array(X)
T_py = np.array(T)
Y_py = np.array(Y)

dml.fit(X_py, T_py, Y_py)
ate = dml.effect(X_py)
\end{minted}

\textbf{Advantages:}
\begin{itemize}
    \item Immediate access to mature Python ecosystem
    \item No reimplementation needed
    \item Gradual migration path
\end{itemize}

\textbf{Disadvantages:}
\begin{itemize}
    \item Python overhead remains
    \item Array copying between languages
    \item Debugging complexity
\end{itemize}

\subsection{Strategy 2: PythonCall.jl (Call Julia from Python)}

Use Julia's performance from Python workflows:

% PythonCall integration
\begin{minted}[fontsize=\small]{python}
# Python
from juliacall import Main as jl

# Load Julia functions
jl.include("dml_julia.jl")

# Call Julia DML (fast!)
ate, se = jl.fit_dml(X, T, Y)
\end{minted}

\textbf{Advantages:}
\begin{itemize}
    \item Keep Python workflow
    \item Julia speed for bottlenecks
    \item Selective optimization
\end{itemize}

\textbf{Disadvantages:}
\begin{itemize}
    \item Julia startup latency (first call)
    \item Maintenance of two codebases
\end{itemize}

\subsection{Strategy 3: Native Julia (Long-term)}

Full reimplementation in Julia:

% Native Julia package structure
\begin{minted}[fontsize=\small]{julia}
module DoubleMachineLearningJulia

using MLJ
using Statistics
using LinearAlgebra

export DoubleML, TemporalPLRDML, PanelDML
export fit!, predict_ate, predict_cate

include("estimators/double_ml.jl")
include("estimators/temporal_plr_dml.jl")
include("estimators/panel_dml.jl")
include("cross_fitting.jl")
include("hac.jl")

end
\end{minted}

\textbf{Advantages:}
\begin{itemize}
    \item Maximum performance
    \item Native Julia ecosystem integration
    \item Type-stable, composable design
\end{itemize}

\textbf{Disadvantages:}
\begin{itemize}
    \item Significant development effort
    \item Validation against established implementations
    \item Maintenance burden
\end{itemize}


\section{Development Roadmap}\label{sec:julia_roadmap}

This section outlines a phased plan for Julia DML implementation.

\subsection{Phase 1: Core Estimators (3-6 months)}

Implement basic DML with MLJ integration:

\begin{itemize}
    \item \textbf{DoubleML}: Standard DML with cross-fitting
    \item \textbf{PartiallyLinearModel}: Robinson estimator
    \item \textbf{InteractiveRegressionModel}: Fully interacted specification
    \item \textbf{TimeSeriesCrossValidator}: Temporal blocking
    \item HAC standard errors with Newey-West
\end{itemize}

Deliverables:
\begin{itemize}
    \item Registered Julia package (General registry)
    \item MLJ integration (model wrappers)
    \item Comprehensive test suite
    \item Documentation with examples
\end{itemize}

\subsection{Phase 2: Extensions (6-12 months)}

Add advanced functionality:

\begin{itemize}
    \item \textbf{TemporalPLRDML}: Scalar temporal PLR effects with HAC inference
    \item \textbf{PanelDML}: Fixed effects with DML
    \item \textbf{CausalForestDML}: Heterogeneity via forest
    \item \textbf{DRIV}: Doubly robust instrumental variables
    \item Stationarity tests (ADF, KPSS, Phillips-Perron)
\end{itemize}

\subsection{Phase 3: Deployment Hardening (12-18 months)}

Deployment hardening requirements:

\begin{itemize}
    \item Model serialization (JLD2.jl)
    \item Causal monitoring system
    \item Distributed computing (Distributed.jl)
    \item GPU acceleration (CUDA.jl)
    \item REST API serving (Genie.jl/HTTP.jl)
\end{itemize}

\subsection{Validation Strategy}

All implementations validated against:

\begin{enumerate}
    \item EconML (Python): Numerical equivalence
    \item Published results: Academic paper replication
    \item Monte Carlo: Coverage and bias verification
    \item This book's test suite: 700+ adapted tests
\end{enumerate}


\section{Code Translation Reference}\label{sec:code_reference}

Quick reference for translating this book's Python code to Julia.

\subsection{Common Operations}

\begin{table}[htbp]
\centering
\caption{Python to Julia operation mapping}
\label{tab:operation_mapping}
\begin{tabular}{ll}
\toprule
\textbf{Python} & \textbf{Julia} \\
\midrule
\texttt{np.mean(x)} & \texttt{mean(x)} \\
\texttt{np.std(x, ddof=1)} & \texttt{std(x)} \\
\texttt{np.clip(x, a, b)} & \texttt{clamp.(x, a, b)} \\
\texttt{np.random.randn(n)} & \texttt{randn(n)} \\
\texttt{x @ y} & \texttt{x * y} \\
\texttt{x * y} (elementwise) & \texttt{x .* y} \\
\texttt{X[mask, :]} & \texttt{X[mask, :]} \\
\texttt{for i in range(n):} & \texttt{for i in 1:n} \\
\texttt{lambda x: x**2} & \texttt{x -> x\textasciicircum 2} \\
\texttt{list comprehension} & \texttt{comprehension} (same syntax) \\
\bottomrule
\end{tabular}
\end{table}

\subsection{Type Annotations}

% Julia type annotations
\begin{minted}[fontsize=\small]{julia}
# Function with type annotations
function fit_dml(
    X::Matrix{Float64},
    T::Vector{Int},
    Y::Vector{Float64};
    n_folds::Int=5
)::DMLResult
    # Implementation
end

# Struct with types
@kwdef struct DMLResult
    ate::Float64
    se::Float64
    ci_lower::Float64 = ate - 1.96 * se
    ci_upper::Float64 = ate + 1.96 * se
end
\end{minted}


\section{Exercises}\label{sec:app_exercises}

\begin{exercise}[Basic Translation]
Translate the HAC standard error computation from Chapter 6 (\texttt{dml\_ts/dml/hac.py}) to Julia. Verify numerical equivalence with the Python implementation on test data.
\end{exercise}

\begin{exercise}[MLJ Integration]
Create an MLJ-compatible wrapper for the DoubleML estimator that integrates with MLJ's \texttt{fit!}, \texttt{predict}, and \texttt{evaluate!} interface. Test with MLJ's built-in models.
\end{exercise}

\begin{exercise}[Performance Profiling]
Using Julia's \texttt{Profile} module and \texttt{BenchmarkTools}, identify the computational bottlenecks in DML cross-fitting. Compare with Python profiling results.
\end{exercise}

\begin{exercise}[PyCall Integration]
Create a Julia script that:
\begin{enumerate}
    \item Loads data in Julia (faster I/O)
    \item Calls EconML's DML via PyCall
    \item Post-processes results in Julia (HAC errors, visualization)
\end{enumerate}
Measure the overhead of Python interop.
\end{exercise}


\section{Summary}\label{sec:app_summary}

This appendix provided a roadmap for implementing causal inference in Julia:

\keyresult{
\textbf{Key Takeaways:}
\begin{itemize}
    \item Julia offers 5-9$\times$ speedup for Monte Carlo/bootstrap workloads
    \item Current Julia ecosystem lacks mature DML (opportunity)
    \item Three integration strategies: PyCall, PythonCall, native
    \item Translation patterns map Python DML to Julia idioms
    \item Phased roadmap: Core (6mo) $\rightarrow$ Extensions (12mo) $\rightarrow$ Production (18mo)
\end{itemize}
}

\vspace{1em}

\begin{tcolorbox}[colback=gray!5, colframe=gray!75, title={Appendix Checklist}]
\begin{itemize}
    \item[$\square$] Understand Julia's performance advantages
    \item[$\square$] Survey Julia ecosystem (MLJ.jl, CausalInference.jl)
    \item[$\square$] Translate Python DML patterns to Julia
    \item[$\square$] Evaluate integration strategies for your workflow
    \item[$\square$] Plan native Julia implementation timeline
\end{itemize}
\end{tcolorbox}

The Julia ecosystem for causal inference is an active area of development. Contributions to packages like CausalInference.jl and new DML implementations will benefit the broader community of researchers and practitioners.
